\documentclass[10pt,a4paper]{amsart}
\usepackage[margin=19mm]{geometry}
\usepackage{amsmath,amssymb,mathtools}
\usepackage[scr=rsfs,cal=euler]{mathalfa}
\usepackage{xfrac}
\usepackage[unicode,hidelinks,bookmarks=false]{hyperref}
\newcommand{\ie}{i.e.\ }
\newcommand{\cf}{cf.\ }
\newcommand{\resp}{respectively\ }
\newcommand{\Subsection}{\subsection}
\providecommand{\nobreakdash}{\nobreak\hbox{-}}
\setlength{\emergencystretch}{2em}
\multlinegap0pt
\usepackage{bm}
\usepackage{bbm}
\usepackage{dsfont}
\usepackage{xspace}
\usepackage{cancel}
\usepackage{mathtools}
\usepackage{amsthm}
\makeatletter
\@ifundefined{theorem}{\newtheorem{theorem}{Theorem}}{}
\@ifundefined{lemma}{\newtheorem{lemma}[theorem]{Lemma}}{}
\makeatother
\newcommand{\E}{ \mathbb{E}}
\renewcommand{\P}{\mathbb{P}}
\title[Mallows Model with Learned Distance Metrics: Sampling and Ma]{Mallows Model with Learned Distance Metrics: Sampling and Maximum Likelihood Estimation}
\author{Yeganeh Alimohammadi, Kiana Asgari}
\date{Source: arXiv:2507.08108v1 (CC BY 4.0)}
\begin{document}
\maketitle

\noindent\emph{Source and licence.} Mallows Model with Learned Distance Metrics: Sampling and Maximum Likelihood Estimation. Version arXiv:2507.08108v1 (CC BY 4.0), distributed under the Creative Commons Attribution 4.0 International licence, as recorded in the arXiv licence metadata for this exact version and shown on the abstract page https://arxiv.org/abs/2507.08108v1. Full rights notice: licenses/arxiv-2507.08108-CC-BY-4.0.txt.\par\smallskip

\noindent\textit{Editorial scope.} Counted unit: the Lemma whose source label is \texttt{lm: identifiable} in the article (source0.tex lines 1309--1317, source label \texttt{lm: identifiable}), delivered together with the article's own complete original proof of it (source0.tex lines 1319--1490, one \texttt{proof} environment of 172 lines). No mathematics is written, repaired or re-derived: statement and proof are the article's own text line for line. Required context retained uncounted: none. Every definition, display and label of those lines is retained unchanged. Omitted from this package and not redistributed: 1--1308, 1318--1318, 1491--2423. Nothing inside a retained excerpt is omitted or altered: every retained body is delivered exactly as the article's lines are written, so no omission receipt of any kind accompanies this package. Citations: the retained excerpts carry no citation command at all, so no bibliography entry of the article is transcribed and no reference is redistributed. Reference adaptation: none was needed and none was made; every reference inside the delivered unit resolves inside it, so no unresolved reference or citation is produced. Printed slips kept literal: no printed word, symbol, hypothesis or number of the counted statement or of its proof was altered. Layout and class: the article's own class is \texttt{article} with packages including arxiv, inputenc, fontenc, hyperref, url, booktabs, natbib, amsfonts, nicefrac, microtype, lipsum, graphicx, amsmath, amssymb; the wrapper is the lane's pinned \texttt{amsart} editorial wrapper at 10pt on A4, which restates the theorem-like environments the retained text needs and adds the article's own macro definitions that the retained text needs (\texttt{\textbackslash E}, \texttt{\textbackslash P}), with extra wrapper packages (bm, bbm, dsfont, xspace, cancel, mathtools). The delivered numbering is the unit's own. Assets: no retained span contains a figure environment, a graphics-inclusion command, a PDF-page inclusion, a table or a TikZ picture; no image file is copied into this package, the package declares no asset dependency and no image file is redistributed with it. Licence: version arXiv:2507.08108v1 of this article is distributed under the Creative Commons Attribution 4.0 International licence, as recorded in the arXiv licence metadata for this exact version and shown on the abstract page https://arxiv.org/abs/2507.08108v1. A full rights notice is delivered at licenses/arxiv-2507.08108-CC-BY-4.0.txt. Characters: every retained excerpt is pure ASCII, so the delivered engine, XeLaTeX with its default Latin Modern text font, reports no missing character.
\begin{lemma}(Identifiability)\label{lm: identifiable}
The population score \(\Psi(\theta)\) satisfies
\(\Psi(\theta_0)=0\) and for every \(\varepsilon>0\),
\[
\inf_{\|\theta-\theta_0\|\ge\varepsilon}
\|\Psi(\theta)\|
>0.
\]
\end{lemma}

\begin{proof}
A straightforward calculation gives
\[
\Psi(\theta)
=
\begin{pmatrix}
\mathbb{E}_{\alpha_0,\beta_0}\left[\dot{d}_\alpha(\pi)\right] - \mathbb{E}_{\alpha,\beta}\left[\dot{d}_\alpha(\pi)\right]
\\[6pt]
\mathbb{E}_{\alpha_0,\beta_0}\left[d_\alpha(\pi)\right] - \mathbb{E}_{\alpha,\beta}\left[d_\alpha(\pi)\right]
\end{pmatrix},
\]
here since $\hat\sigma_m$ is fixed, with abuse of notation we write $d_\alpha(\pi)$ instead of $d_\alpha(\pi,\hat\sigma_m)$ and $\E_{\alpha',\beta'}$ instead of $\E_{\P{_\alpha',\beta',\hat \sigma_m}}$.
With this, proving identifiability reduces to showing that the system of equations
\[
\mathbb{E}_{\alpha,\beta}[d_\alpha(\pi)] = \mathbb{E}_{\alpha_0,\beta_0}[d_\alpha(\pi)],
\quad
\mathbb{E}_{\alpha,\beta}[\dot{d}_\alpha(\pi)] = \mathbb{E}_{\alpha_0,\beta_0}[\dot{d}_\alpha(\pi)]
\]
admits a unique solution \((\alpha,\beta) = (\alpha_0,\beta_0)\).

To facilitate the argument, we fix \((\alpha,\beta)\) and view the left-hand sides as functions of \((\alpha_0,\beta_0)\). In this formulation, identifiability reduces to proving that for each fixed \((\alpha,\beta)\), there exists a unique \((\alpha_0,\beta_0)\) satisfying the system.

For this purpose, we introduce the auxiliary map
\[
\boldsymbol{f}_{\alpha}(\alpha_0,\beta_0)
=
\left(
\mathbb{E}_{{\alpha_0,\beta_0}}[d_{\alpha}(\pi)],\;
\mathbb{E}_{{\alpha_0,\beta_0}}[\dot{d}_{\alpha}(\pi,)]
\right),
\]
where \(\alpha \) is fixed. That is, while the distance function $d_\alpha$ is fixed, the true data-generating parameters \((\alpha_0,\beta_0)\) varies.
The key idea is that if \(\boldsymbol{f}_{\alpha}\) is injective, then the solution to the system is unique. We now proceed to prove that \(\boldsymbol{f}_{\alpha}\) is a global diffeomorphism onto its image. To this end, we first establish local invertibility and properness, which together imply global invertibility by the Hadamard--Caccioppoli theorem.



\begin{lemma}[Local Invertibility]\label{lem:local_invertibility}
For any \(\alpha > 0\), the function \(\boldsymbol{f}_{\alpha}\) is locally invertible on \((0,\infty)^2\).
\end{lemma}

\begin{proof}
It suffices to show that the Jacobian matrix \(\boldsymbol{J}\boldsymbol{f}_{\alpha}(\alpha_0,\beta_0)\) is invertible at every point \((\alpha_0,\beta_0) \in (0,\infty)^2\).
Direct calculation gives:
\[
\boldsymbol{J}\boldsymbol{f}_{\alpha}(\alpha_0,\beta_0)
= -
\begin{bmatrix}
\mathrm{Cov}_{\alpha_0,\beta_0}\left(d_\alpha(\pi), d_{\alpha_0}(\pi)\right) &
\beta \, \mathrm{Cov}_{\alpha_0,\beta_0}\left(d_\alpha(\pi), \dot{d}_{\alpha_0}(\pi)\right)
\\[6pt]
\mathrm{Cov}_{\alpha_0,\beta_0}\left(\dot{d}_\alpha(\pi), d_{\alpha_0}(\pi)\right) &
\beta \, \mathrm{Cov}_{\alpha_0,\beta_0}\left(\dot{d}_\alpha(\pi), \dot{d}_{\alpha_0}(\pi)\right)
\end{bmatrix}.
\]

Define the two-dimensional random vectors:
\[
\boldsymbol{a}_1(\pi) = \left(d_{\alpha_0}(\pi), \dot{d}_{\alpha_0}(\pi)\right),
\quad
\boldsymbol{a}_2(\pi) = \left(d_{\alpha}(\pi), \dot{d}_{\alpha}(\pi)\right).
\]
Then
\[
\det\left( \boldsymbol{J}\boldsymbol{f}_{\tilde\alpha}(\alpha,\beta) \right)
= \beta \cdot \det\left( \mathrm{Cov}_{\alpha,\beta}\left(\boldsymbol{a}_1(\pi), \boldsymbol{a}_2(\pi)\right) \right).
\]

When \((\alpha_0,\beta_0) = (\alpha,\beta)\), we have \(\boldsymbol{a}_1(\pi) = \boldsymbol{a}_2(\pi)\), and the covariance matrix reduces to the variance of \(\boldsymbol{a}_1(\pi)\), which is positive definite.

When \(\alpha_0 \neq \alpha\),  
suppose, for contradiction, that \(\mathrm{Cov}_{\alpha_0,\beta_0}(\boldsymbol{a}_1(\Pi), \boldsymbol{a}_2(\Pi))\) is singular. Then there must exist nonzero vectors \(\lambda^{(1)}, \lambda^{(2)} \in \mathbb{R}^2\) such that
\begin{equation}\label{eq:invertibility_f_jacobian}
\left\langle \lambda^{(1)}, \boldsymbol{a}_1(\Pi) \right\rangle
+
\left\langle \lambda^{(2)}, \boldsymbol{a}_2(\Pi) \right\rangle
\quad \text{is constant almost surely}.
\end{equation}
Since \(\P_{\alpha_0,\beta_0}\) has full support on the finite set \(S_n\), this equality holds for every \(\pi \in S_n\).


To derive a contradiction, we consider a few examples. Let \(\pi_1=(2,4)\) denote the transposition swapping elements \(2\) and \(4\), and \(\pi_2\) denote the composition of \(\pi_1\) with the transposition swapping \(1\) and \(3\), that is, \(\pi_2 = (1\;3)(2\;4)\).
Applying \eqref{eq:invertibility_f_jacobian} to both \(\pi_1\) and \(\pi_2\), we have:
\begin{align*}
\left\langle \lambda^{(1)}, \boldsymbol{a}_1(\pi_1) \right\rangle + \left\langle \lambda^{(2)}, \boldsymbol{a}_2(\pi_1) \right\rangle &= 
\left\langle \lambda^{(1)}, \boldsymbol{a}_1(\pi_2) \right\rangle + \left\langle \lambda^{(2)}, \boldsymbol{a}_2(\pi_2) \right\rangle.
\end{align*}
By construction, we observe:
\[
\boldsymbol{a}_1(\pi_2) = \boldsymbol{a}_1(\pi_1) + \boldsymbol{a}_1((1\;3)),
\quad
\boldsymbol{a}_2(\pi_2) = \boldsymbol{a}_2(\pi_1) + \boldsymbol{a}_2((1\;3)).
\]
Substituting the expressions for \(\pi_2\), we obtain:
\[
\left\langle \lambda^{(1)}, \boldsymbol{a}_1((1\;3)) \right\rangle
+
\left\langle \lambda^{(2)}, \boldsymbol{a}_2((1\;3)) \right\rangle
= 0.
\]

Repeating this argument with other simple transpositions such as \((1\;4)\), \((1\;5)\), and \((1\;6)\), we obtain a homogeneous linear system:
\[
\boldsymbol{A}
\begin{pmatrix}
\lambda^{(1)} \\ \lambda^{(2)}
\end{pmatrix}
= 0,
\]
where
\[
\boldsymbol{A}
=
\begin{bmatrix}
\boldsymbol{a}_1(1\;3) & \boldsymbol{a}_2(1\;3) \\
\boldsymbol{a}_1(1\;4) & \boldsymbol{a}_2(1\;4) \\
\boldsymbol{a}_1(1\;5) & \boldsymbol{a}_2(1\;5) \\
\boldsymbol{a}_1(1\;6) & \boldsymbol{a}_2(1\;6)
\end{bmatrix}.
\]
It can be checked that when \(\alpha \neq \alpha_0\), the matrix \(\boldsymbol{A}\) has full rank, 
Since the points \((d_{\alpha_0}(\pi), \dot{d}_{\alpha_0}(\pi), d_{\alpha}(\pi), \dot{d}_{\alpha}(\pi))\) for different transpositions \((1\;k)\) involve distinct powers of \(|k-1|\) under \(\alpha\) and \(\alpha_0\), and the functions \(\alpha \mapsto |k-1|^\alpha\) and \(\alpha \mapsto \log|k-1|\cdot |k-1|^\alpha\) are linearly independent for different \(k\), the matrix \(\boldsymbol{A}\) is full rank whenever \(\alpha \neq \alpha_0\).
As a result, \(\lambda^{(1)} = \lambda^{(2)} = 0\).
Thus, no nontrivial linear relation exists, contradicting the assumption that \(\mathrm{Cov}_{\alpha_0,\beta_0}(\boldsymbol{a}_1(\Pi), \boldsymbol{a}_2(\Pi))\) is singular. This completes the proof of local invertibility.
\end{proof}

\bigskip

\begin{lemma}[Properness]\label{lem:properness}
For any \(\alpha\), the function \(\boldsymbol{f}_{\alpha}\) is proper: the preimage of every compact set is compact.
\end{lemma}

\begin{proof}
Let \(K\subset \mathrm{Img}(\boldsymbol{f}_{\tilde\alpha})\) be any compact set.  
We must show that \(\boldsymbol{f}_{\tilde\alpha}^{-1}(K)\) is compact.
Suppose for contradiction that \(\boldsymbol{f}_{\tilde\alpha}^{-1}(K)\) is not compact.  
Then there exists a sequence \((\alpha_m,\beta_m)\) in \(\boldsymbol{f}_{\tilde\alpha}^{-1}(K)\) that escapes to infinity, meaning
\[
\lim_{m\to\infty} \|(\alpha_m,\beta_m)\| = \infty.
\]

As \((\alpha_m,\beta_m)\to\infty\), the corresponding distributions \(\P_{\alpha_m,\beta_m}\) concentrate on the identity permutation. To see this as \((\alpha_m,\beta_m) \to \infty\), either \(\alpha_m \to \infty\) or \(\beta_m \to \infty\) or both. In any of these cases, for any permutation \(\pi \neq \mathrm{id}\), we have
\[
\frac{\P_{\alpha_m,\beta_m}(\pi)}{\P_{\alpha_m,\beta_m}(\mathrm{id})} = \exp\left( -\beta_m ( d_{\alpha_m}(\pi) - d_{\alpha_m}(\mathrm{id}) ) \right) \to 0.
\] Thus, as $m$ increases,
\[
\boldsymbol{f}_{\tilde\alpha}(\alpha_m, \beta_m) \to \boldsymbol{0}.
\]
But note that for any finite \((\alpha,\beta) \in (0,\infty)^2\), we always have:
\[
\mathbb{E}_{\alpha,\beta}[d_{\tilde\alpha}(\Pi)] > 0,
\quad
\mathbb{E}_{\alpha,\beta}[\dot{d}_{\tilde\alpha}(\Pi)] > 0,
\]
since \(\P_{\alpha,\beta}\) has full support on \(S_n\), and \(d_{\tilde\alpha}(\Pi) > 0\) with positive probability.
Therefore, \(\boldsymbol{0} \notin \mathrm{Img}(\boldsymbol{f}_{\tilde\alpha})\).
Now, since \(\boldsymbol{f}_{\tilde\alpha}(\alpha_m, \beta_m) \in K\) for all \(m\), and \(\boldsymbol{f}_{\tilde\alpha}(\alpha_m, \beta_m) \to \boldsymbol{0}\), we conclude that \(\boldsymbol{0} \in K\) by closedness of \(K\).
This contradicts \(\boldsymbol{0} \notin \mathrm{Img}(\boldsymbol{f}_{\tilde\alpha})\). Therefore, the preimage \(\boldsymbol{f}_{\tilde\alpha}^{-1}(K)\) must be compact.

\end{proof}

Now, to finish the proof of uniqueness of the critical point, observe that the condition \(\Psi(\theta) = 0\) is equivalent to
\[
\boldsymbol{f}_\alpha(\alpha,\beta) = \boldsymbol{f}_\alpha(\alpha_0,\beta_0).
\]
Since \(\boldsymbol{f}_\alpha\) is a global diffeomorphism (by Lemmas~\ref{lem:local_invertibility} and~\ref{lem:properness}), it follows that \((\alpha,\beta) = (\alpha_0,\beta_0)\), and hence \(\Psi(\theta_0) = 0\) has a unique solution.

To complete the identifiability condition, we note that the global invertibility of \(\boldsymbol{f}_\alpha\) implies that its inverse \(\boldsymbol{f}_\alpha^{-1}\) is continuous on its image. Therefore, \(\Psi(\theta) = \boldsymbol{f}_\alpha(\alpha_0, \beta_0) - \boldsymbol{f}_\alpha(\theta)\) is continuous, vanishes only at \(\theta_0\), and cannot tend to zero along any sequence staying at distance at least \(\varepsilon > 0\) from \(\theta_0\). Thus, for every \(\varepsilon > 0\),
\[
\inf_{\|\theta - \theta_0\| \ge \varepsilon} \|\Psi(\theta)\| > 0,
\]
as required.
\end{proof}

\end{document}
